\ssr{ЗАКЛЮЧЕНИЕ}

В результате выполнения практики была достигнута поставленная цель -- разработана модификация муравьиного алгоритма с динамическим подбором гетерогенных параметров на основе генетического алгоритма для решения задачи коммивояжёра.

Было описано и проведено функциональное тестирование. Во всех запусках сформирован гамильтонов цикл, алгоритм завершал работу без ошибок, параметры $\alpha$ и $\beta$ менялись индивидуально у каждого муравья, все функциональные тесты пройдены.

Было проведено модульное тестирование с использованием testing и testify, покрыты модули стратегий, структур данных и части ядра, покрытие кода составило 42,2\%.

Было проведено исследование эволюционных операторов. Наиболее удачной конфигурацией оказалось сочетание рулеточной селекции, BLX-скрещивания и гауссовской мутации.

Было проведено исследование сходимости. Разработанный алгоритм без локальной оптимизации достигает качества классического муравьиного алгоритма (разница 1–2\%) в 3.7-8 раз быстрее по числу итераций. Добавление 2-opt улучшает качество решений на 6–8\%.

Было проведено исследование устойчивости к параметрам $\alpha$ и $\beta$. Разработанная модификация показала коэффициент вариации средней длины маршрута 4,54\% против 27,96\% у классического алгоритма.